\clearpage
\begin{trapbox}
\textbf{先看反例：方向相同，也不能保证弦长相同。}
\begin{center}\begin{tikzpicture}[font=\small,>=stealth]
\draw[rounded corners=3pt,gray!45] (0,0) rectangle (14.2,10.8);
\node at (7.1,10.25) {先检查直线：轴线只有一个交点；竖直弦也必须经过焦点};
\node at (2.6,9.3) {轴线：不构成焦点弦};
\begin{scope}[shift={(2,5.3)},scale=.68]
\draw[->,gray] (-1.5,0)--(4.4,0) node[above] {$x$};
\draw[teal!75!black,thick,samples=120,domain=-3.7:3.7,variable=\t] plot ({\t*\t/4},{\t});
\draw[red!75!black,very thick] (-.6,0)--(4.1,0);\fill (0,0) circle(2pt) node[below left=5pt] {$O$};
\fill (1,0) circle(2pt) node[below right=5pt] {$F$};
\end{scope}
\node at (2.6,1.85) {$\theta=0$，仅顶点 $O$};
\node at (7.0,9.3) {通径：合法且最短};
\begin{scope}[shift={(6.4,5.3)},scale=.68]
\draw[->,gray] (-1.5,0)--(4.4,0) node[above] {$x$};
\draw[teal!75!black,thick,samples=120,domain=-3.7:3.7,variable=\t] plot ({\t*\t/4},{\t});
\draw[blue!75!black,very thick] (1,-2.0)--(1,2.0);
\fill (1,2.0) circle(2pt);\fill (1,-2.0) circle(2pt);
\fill (1,0) circle(2pt) node[below right=5pt] {$F$};
\end{scope}
\node at (7.0,1.85) {$x=1,\quad L=4$};
\node at (11.4,9.3) {普通弦：不能套公式};
\begin{scope}[shift={(10.8,5.3)},scale=.68]
\draw[->,gray] (-1.5,0)--(4.4,0) node[above] {$x$};
\draw[teal!75!black,thick,samples=120,domain=-3.7:3.7,variable=\t] plot ({\t*\t/4},{\t});
\draw[blue!75!black,very thick] (2,-2.8284271247461903)--(2,2.8284271247461903);
\fill (2,2.8284271247461903) circle(2pt);\fill (2,-2.8284271247461903) circle(2pt);
\fill (1,0) circle(2pt) node[above=8pt] {$F$};
\end{scope}
\node at (11.4,1.85) {$x=2,\quad L=4\sqrt2$};

\node at (7.1,.65) {三图均为 $y^2=4x$，同尺度；后两条直线同为竖直，只有中图过 $F$。};
\end{tikzpicture}
\end{center}
左图轴线仅交于顶点，不能当作两端点不同的焦点弦；中、右两图同为竖直弦，右图不经过焦点，所以长度不同。
\textbf{正确做法：}先验“过焦点、两个不同交点”，再套弦长公式。
当 $\theta\to0^+$ 或 $\pi^-$ 时弦长趋于无穷大，是极限现象；$\theta=0$ 本身没有这条有限焦点弦。
\end{trapbox}
\clearpage
\begin{trapbox}
\textbf{开口方向变了，先看角从哪条轴量。}
\begin{center}\begin{tikzpicture}[font=\small,>=stealth]
\draw[rounded corners=3pt,gray!45] (0,0) rectangle (14.2,12);
\node at (7.1,11.45) {开口向上时：角仍从横轴量，分母就要改用 $\cos^2\theta$};

\begin{scope}[shift={(4.4,3.3)},scale=1.05]
\draw[->,gray] (-2.4,0)--(4.3,0) node[above] {$x$};
\draw[->,gray] (0,-.4)--(0,4.8) node[left] {$y$};
\draw[teal!75!black,thick,samples=120,domain=-2.2:4.1,variable=\t] plot ({\t},{\t*\t/4});
\draw[blue!75!black,very thick] (0,1)--(3.464101615,3);
\draw[orange!85!black,very thick] (0,1)--(-1.154700538,.333333333);
\draw[gray,dashed] (0,1)--(2,1);
\draw[violet,->] (.8,1) arc[start angle=0,end angle=30,radius=.8];
\node[violet] at (1.3,1.28) {$\theta$};
\draw[orange!85!black,->] (.95262794,1.55) arc[start angle=30,end angle=90,radius=1.1];
\node[orange!85!black] at (.62,2.34) {$\varphi$};
\fill (0,1) circle(2pt) node[above left=5pt] {$F$};
\fill (3.464101615,3) circle(2pt) node[above left=5pt] {$A$};
\fill (-1.154700538,.333333333) circle(2pt) node[above=10pt] {$B$};
\node[below left=4pt] at (0,0) {$O$};
\end{scope}
\node[anchor=west,align=left] at (9.4,7)
 {$x^2=4y$，$F=(0,1)$\\[12pt]
 $\textcolor{violet}{\theta=30^\circ}$（从横轴量）\\[9pt]
 $\textcolor{orange!85!black}{\varphi=60^\circ}$（与竖轴）\\[12pt]
 $\sin^2\varphi=\cos^2\theta$\\[9pt]
 $L=4/\cos^230^\circ=16/3$};
\node at (7.1,1.5) {先认对称轴，再认角度：图中 $\theta+\varphi=90^\circ$。};
\node at (7.1,.65) {通用记法：用弦与对称轴的夹角 $\varphi$，总有 $L=2p/\sin^2\varphi$。};
\end{tikzpicture}
\end{center}
左右开口、角从 $x$ 轴量时用 $\sin^2\theta$；上下开口、角仍从 $x$ 轴量时改用 $\cos^2\theta$。
图中的弦长是 $16/3$，不能把 $30^\circ$ 直接代进分母 $\sin^2\theta$。
最短的始终是垂直于对称轴的通径：上下开口时为水平弦。
\end{trapbox}
\clearpage
\begin{trapbox}
\begin{description}[style=nextline,leftmargin=2.5em,itemsep=.5em]
\item[\textbf{易错一：把方程系数和 $p$ 混为一谈。}]
$y^2=8x$ 表示 $2p=8$，\textbf{$p=4$}；通径长度 $2p=8$，不是 $16$。若写 $y^2=4ax$，则本结论中 $p=2a$、通径长 $4a$。
\item[\textbf{易错二：给左开口硬套负号。}]
$y^2=-2px$ 里的 $p>0$，弦长是\textbf{正长度}：$L=2p/\sin^2\theta$，不是负的。
\item[\textbf{易错三：用斜率 $k$ 冒充 $\sin\theta$。}]
$\sin^2\theta=k^2/(1+k^2)$（有限 $k\ne0$），故斜率式是
$L=2p(1+k^2)/k^2$，\textbf{不能写 $2p/k^2$}。
\item[\textbf{易错四：漏掉竖直弦。}]
竖直弦 $\theta=\pi/2$，斜率不存在，却\textbf{完全合法且最短}；直接取 $L=2p$。
\item[\textbf{易错五：误把轴线看成一条焦点弦。}]
水平轴线方向 $\theta=0$（$\pi$ 也非标准倾斜角取值），过焦点的轴线\textbf{仅在顶点}与抛物线相交，不形成两端点不同的有限焦点弦；不能将 $\sin\theta=0$ 当作无穷大的实际弦长。只有当 $\theta\to0^+$ 或 $\pi^-$ 时，弦长才趋于 $+\infty$。
\item[\textbf{易错六：认为角度越大弦越短。}]
最短是在 $90^\circ$，而 $30^\circ$ 与 $150^\circ$ 的弦长相等；真正规律是\textbf{越靠近 $90^\circ$ 越短}。
\item[\textbf{易错七：把上下开口仍用 $\sin^2\theta$。}]
$x^2=\pm2py$ 的轴线竖直；倾角从 $x$ 轴量时用
$L=2p/\cos^2\theta$，最短焦点弦为\textbf{水平}，$\theta=0$。
\item[\textbf{易错八：把“过焦点”遗漏。}]
$L=2p/\sin^2\theta$ \textbf{只用于过焦点的弦}；一般平行弦长度不同，不能只靠角度算。
\item[\textbf{易错九：焦半径和等于弦长的前提。}]
本题 $A,F,B$ 共线且焦点\textbf{处于两端点之间}，才有 $|AB|=|AF|+|BF|$；其他三点排列可能需要作差。
\end{description}
\medskip
\textbf{\textcolor{CTrap}{最稳的三项核对：}}$L\ge 2p$？左右开口用的都是正 $p$？竖直弦得的是 $2p$？
\end{trapbox}
